\scanpage{4}\mspage{3}
À cause de cette interprétation, le passage de $\widetilde F\subset M$ à $\operatorname{Fig}_{M}(F)$ est à tel point évident, que \inserted{la \unsure{considération} de} cette figure ensembliste dans $M$ me semble \unsure{rarement|vraiment} importante -- mais à voir.....

\deleted{Quand on identi}

On pourra utiliser les deux applications (1.3) pour identifier une figure, suivant les cas, \unsure{à} une partie de $M$ [via $F\mapstochar\relbar\widetilde F$] ou à une figure ensembliste dans $M$ [via $F\mapsto\operatorname{\text{Dépl}}(F)$]. Dans le premier cas, \unsure{comparant|composant} \unsure{avec|à un} $M\unsure{\lhook\joinrel\relbar|\hookrightarrow}\mathcal{F}$ (1.2)\deleted{\unsure{(}}, par la formule 1,5. Cette dernière inclusion s'identifie à
\[
X\mathrel{\mapstochar\relbar\mapstochar}\widetilde X=\{Y\in M\mid Y\leq X\}
\qquad M\to\mathcal{P}_{\mathrm{ferm}}(M).
\]

Donc il faudra être prudent alors \inserted{\unsure{avec}} l'identific\unsure{ation} dont il a été question, des\deleted{\unsure{une|un}} $\unsure{X|X'}\in M$ \unsure{ou|en} \unsure{multistrates|multistrate} à des figures particulières, quand celles-ci sont interprétées comme des parties de $M$: La figure associée à $X\in M$ n'est nullement $\{X\}$, mais $\widetilde X\rbrace$ \deleted{\unsure{$f$|$\mathcal F$|$F$|l'}} ens des multistrates majorées par $x$ pour $\leq$ (ou encore, des sous-multistrates, ou multistrates incidentes à $X$, ou encore des \deleted{\unsure{$f$|$\mathcal F$|$F$}} \uline{strates} de $x$).
